% Generated by tools/edn_latex.py from reports/edn.md.
% Do not edit: change reports/edn.md and run `make edn`.
\ednentry{
  id = {EDN-16},
  title = {UC1 and UC2 have separate endpoints},
  date = {2026-09-29},
  milestone = {M1: Project Inception},
  topic = {API Design},
  participants = {@lukas2510},
  decision = {UC1 and UC2 keep their own endpoints, \texttt{POST /\allowbreak{}predict} and \texttt{POST /\allowbreak{}price-\allowbreak{}range} (FR-06, FR-07).},
  rationale = {\emph{Alternatives considered:}
\begin{itemize}[leftmargin=*, topsep=2pt]
\item One \texttt{/\allowbreak{}valuation} endpoint returning the estimate and the interval: the required fields would depend on the use case inside one schema, and both kinds of traffic would land in one log.
\end{itemize}
\par
\emph{Why:} The endpoints differ in their input contract: \texttt{/\allowbreak{}predict} requires ten fields (EDN-15) and \texttt{/\allowbreak{}price-\allowbreak{}range} far fewer.
One endpoint would hide two contracts in a conditional schema.
Monitoring settles it: a change in the mix of clients would look like drift on the completeness features the \texttt{ES} scenario relies on (EDN-14).
A combined endpoint can be added later without breaking clients.},
  ai = {Alternative assessment, Recommendation},
  response = {Accepted},
  assessment = {AI contributed the monitoring argument, which was written down nowhere, and withdrew its own first claim that the question blocked the merge.
Lukas dropped its further suggestions, a shared base schema and a combined convenience endpoint, as unnecessary for a course project.},
  aievidence = {Claude Code session of 2026-09-29; Lukas replied (translated from German): ``I do not want overkill, it should stay simple, it is only a university course''.},
  evidence = {\href{https://github.com/mlops-2627q1-mds-upc/MLOPS_Recommenditos/blob/main/docs/docs/specification.md}{Specification} FR-06 and FR-07, \href{https://github.com/mlops-2627q1-mds-upc/MLOPS_Recommenditos/pull/19}{PR \#19}.},
}
